% md2pdf-preamble.tex — Personalizacións LaTeX para md2pdf
\usepackage{colortbl}
\providecommand{\pandocbounded}[1]{#1}
% ── Cores configurables ───────────────────────────────────────
\newcommand{\toctitlecolor}{blue!45!black}
\newcommand{\titlebodycolor}{heading-section}
\newcommand{\tablestripecolor}{gray!6!white}
\newcommand{\tablebasecolor}{white}
% ── Cabeceira, pé e lendas de táboa ──────────────────────────
\KOMAoptions{headsepline=false,footsepline=false}
\AfterEndPreamble{%
  \KOMAoptions{headsepline=false,footsepline=false}%
  \clearpairofpagestyles
  \cfoot{\pagemark}%
  \pagestyle{scrheadings}%
}
% ── Aliñamento á esquerda ───────────────────────────────────
\usepackage{ragged2e}
\setlength{\RaggedRightRightskip}{0pt plus 2.5em}
\AtBeginDocument{\RaggedRight}
% ── Microtipografía ─────────────────────────────────────────
\AtBeginDocument{%
  \microtypesetup{activate={true,nocompatibility}}%
}

% ── Orfas, viúvas e paxinación ────────────────────────────────
\widowpenalty=3000
\clubpenalty=3000
\displaywidowpenalty=2000

% ── Hifenación e saltos de páxina ─────────────────────────────
\brokenpenalty=5000
\doublehyphendemerits=10000
\finalhyphendemerits=5000

% ── Notas ao pé ────────────────────────────────────────────────
\interfootnotelinepenalty=8000

% ── Protección de cabeceiras ───────────────────────────────
\usepackage{needspace}
% ── Hifenación multilingüe ──────────────────────────────────
\tolerance=1500
\hbadness=1500
\emergencystretch=1.5em
\frenchspacing
\hfuzz=0.5pt
\hyphenation{
  func-tion con-trol-ler mid-dle-ware frame-work data-base name-space
  call-back tem-plate over-ride in-ter-face in-her-it-ance ex-cep-tion
  re-quest re-sponse rou-ting de-bug-ger break-point run-time
  type-script java-script check-box drop-down tool-bar work-flow
  back-end front-end web-hook end-point auto-load auto-com-plete
  over-flow under-score time-stamp fire-wall to-ken ses-sion brow-ser
  ser-ver cli-ent
}
% ── Espazado vertical e cabeceiras ──────────────────────────
\usepackage{setspace}
% Entreliña 1,45 veces o tamaño (13 pt × 1,2 × 1,2083 = 18,85 pt), o límite
% superior do intervalo tipográfico habitual (1,2–1,45): a liña ten uns 80
% caracteres e a Source Serif 4 ten a altura da x alta (0,475 em). As táboas
% herdan o mesmo factor, así que texto e celas teñen a mesma proporción.
\setstretch{1.2083}
\RedeclareSectionCommand[beforeskip=1.6em, afterskip=0.6em]{section}
\RedeclareSectionCommand[beforeskip=1.3em, afterskip=0.4em]{subsection}
\RedeclareSectionCommand[beforeskip=1.0em, afterskip=0.3em]{subsubsection}
\definecolor{heading-section}{HTML}{1C2F4A}
\definecolor{heading-subsection}{HTML}{2D5380}
\definecolor{heading-subsubsection}{HTML}{4A6A96}
\addtokomafont{section}{\color{heading-section}}
\addtokomafont{subsection}{\color{heading-subsection}}
\addtokomafont{subsubsection}{\color{heading-subsubsection}}
\setlength{\intextsep}{8pt plus 2pt minus 2pt}
\setlength{\abovecaptionskip}{3pt}
\setlength{\belowcaptionskip}{1pt}
\AtEndEnvironment{figure}{\vspace{-8pt}}
% Se a cabeceira vén xusto despois doutra (\if@nobreak segue activo),
% non se reserva espazo: \Needspace forzaría un salto entre ambas.
\makeatletter
\newcommand{\mdpdf@headneedspace}[1]{\if@nobreak\else\Needspace{#1}\fi}
\AtBeginDocument{%
  \pretocmd{\section}{\mdpdf@headneedspace{8\baselineskip}}{}{}%
  \pretocmd{\subsection}{\mdpdf@headneedspace{7\baselineskip}}{}{}%
  \pretocmd{\subsubsection}{\mdpdf@headneedspace{6\baselineskip}}{}{}%
}
\makeatother
\AfterEndPreamble{%
  \setlength{\topsep}{0.15em}%
  \setlength{\partopsep}{0pt}%
  \setlength{\itemsep}{0.3em}%
  \setlength{\parsep}{0.2em}%
}
% ── Blockquotes ─────────────────────────────────────────────
\usepackage{tcolorbox}
\tcbuselibrary{breakable,skins}
% Estilo baleiro que callouts.lua aplica ao final de cada callout;
% os snippets --print/--night de lualatex.py redefínenno.
\tcbset{md2pdf-callout/.style={}}
\renewenvironment{quote}{%
  \begin{tcolorbox}[
    colback=blue!3!white, colframe=blue!40!white,
    leftrule=3pt, rightrule=0pt, toprule=0pt, bottomrule=0pt,
    arc=2pt, outer arc=2pt, boxsep=3pt, left=8pt, right=8pt,
    top=4pt, bottom=4pt, before skip=0.45em, after skip=0.6em,
    breakable
  ]%
}{%
  \end{tcolorbox}%
}
% ── Imaxes, figuras e lendas ─────────────────────────────────
\usepackage[export]{adjustbox}
\AtBeginDocument{\setkeys{Gin}{max width=\linewidth}}
\usepackage{caption}
\AtBeginDocument{%
  \captionsetup{format=hang, justification=RaggedRight, font=small, labelfont=bf, labelsep=period, singlelinecheck=true}%
  \captionsetup[table]{position=above}%
}
% ── Código en liña ──────────────────────────────────────────
\definecolor{inline-fg}{HTML}{0550ae}
\let\oldtexttt\texttt
\renewcommand{\texttt}[1]{\textcolor{inline-fg}{\oldtexttt{#1}}}
% ── Til literal ──────────────────────────────────────────────
\AtBeginDocument{%
  \protected\def\textasciitilde{\char126\relax}%
}
% ── Bloques de código ───────────────────────────────────────
\definecolor{code-bg}{HTML}{f6f8fa}
\definecolor{code-border}{HTML}{e1e4e8}
\usepackage{fvextra}
\fvset{breaklines=true, breakanywhere=true}
\DefineVerbatimEnvironment{verbatim}{Verbatim}{breaklines=true,breakanywhere=true}
\makeatletter
\@ifundefined{Shaded}{%
  \newenvironment{Shaded}{}{}%
}{%
  \renewenvironment{Shaded}{}{}%
}%
\renewenvironment{Shaded}{%
  \begin{tcolorbox}[
    colback=code-bg, colframe=code-border, boxrule=0.5pt,
    arc=3pt, outer arc=3pt, fontupper=\footnotesize,
    left=6pt, right=6pt, top=4pt, bottom=4pt,
    before skip=0.45em, after skip=0.7em, breakable,
    lines before break=3
  ]%
}{%
  \end{tcolorbox}%
}
\makeatother
% ── Táboas ───────────────────────────────────────────────────
\makeatletter
\let\md@orig@midrule\midrule
\let\md@orig@toprule\toprule
\let\md@orig@bottomrule\bottomrule
\renewcommand{\midrule}{\noalign{\vspace{2pt}}}
\renewcommand{\toprule}{}
\renewcommand{\bottomrule}{}
\newcommand{\md@top@minipage}{%
  \@ifnextchar[{\md@top@minipage@i}{\md@orig@minipage}%
}
\def\md@top@minipage@i[#1]{\md@orig@minipage[t]}
% arraystretch 1.5: entre fila e fila queda media liña máis ca entre as liñas
% dunha mesma cela (8,57 pt sobre 17,14 pt), a mesma proporción ca \parskip no
% corpo. Con 1.2 o engadido era 0.2 liñas e as filas de varias liñas
% confundíanse coas veciñas.
\AtBeginEnvironment{longtable}{%
  \let\md@orig@minipage\minipage
  \let\minipage\md@top@minipage
  \renewcommand{\arraystretch}{1.5}%
  \small%
  \addfontfeatures{Numbers=Lining}%
  \setlength{\tabcolsep}{6pt}%
  \setlength{\heavyrulewidth}{0.4pt}%
  \arrayrulecolor{heading-section}%
  \sloppy\hbadness=9999\emergencystretch=3em%
  \rowcolors{2}{\tablestripecolor}{\tablebasecolor}%
}
\AtBeginEnvironment{tabular}{%
  \let\md@orig@minipage\minipage
  \let\minipage\md@top@minipage
  \renewcommand{\arraystretch}{1.5}%
  \small%
  \addfontfeatures{Numbers=Lining}%
  \setlength{\tabcolsep}{6pt}%
  \setlength{\heavyrulewidth}{0.4pt}%
  \arrayrulecolor{heading-section}%
  \sloppy\hbadness=9999\emergencystretch=3em%
  \rowcolors{2}{\tablestripecolor}{\tablebasecolor}%
}
\AtBeginEnvironment{tabularx}{%
  \let\md@orig@minipage\minipage
  \let\minipage\md@top@minipage
  \renewcommand{\arraystretch}{1.5}%
  \small%
  \addfontfeatures{Numbers=Lining}%
  \setlength{\tabcolsep}{6pt}%
  \setlength{\heavyrulewidth}{0.4pt}%
  \arrayrulecolor{heading-section}%
  \sloppy\hbadness=9999\emergencystretch=3em%
}
\makeatother
% ── Índice (opción B centrada: caixa restrinxida sen puntos guía) ─────
\AtBeginDocument{%
  \DeclareTOCStyleEntry[
    beforeskip=1.1em plus 1pt,
    entryformat={\bfseries\color{heading-section}},
    pagenumberformat={\bfseries\color{\toctitlecolor}},
    linefill=\hfill
  ]{tocline}{section}%
  \DeclareTOCStyleEntry[
    beforeskip=0.2em,
    entryformat={\color{heading-subsection}},
    pagenumberformat={\color{\toctitlecolor}},
    linefill=\hfill
  ]{tocline}{subsection}%
  \DeclareTOCStyleEntry[
    beforeskip=0.15em,
    entryformat={\color{heading-subsubsection}},
    pagenumberformat={\color{\toctitlecolor}},
    linefill=\hfill
  ]{tocline}{subsubsection}%
  \let\origtableofcontents\tableofcontents
  \renewcommand{\tableofcontents}{%
    \par{\centering
    \begin{minipage}{0.68\linewidth}%
      \origtableofcontents
    \end{minipage}\par}%
  }%
}
% ── Iconas FontAwesome5 ───────────────────────────────────────
\usepackage{fontawesome5}
\DeclareRobustCommand{\iconwarning}{{\color{orange!80!black}\faExclamationTriangle}}
\DeclareRobustCommand{\icontip}{{\color{blue!70!black}\faLightbulb}}
\DeclareRobustCommand{\iconinfo}{{\color{blue!70!black}\faInfoCircle}}
\DeclareRobustCommand{\iconerror}{{\color{red!70!black}\faTimesCircle}}
\DeclareRobustCommand{\iconok}{{\color{green!60!black}\faCheckCircle}}
\DeclareRobustCommand{\iconbug}{{\color{red!60!black}\faBug}}
\DeclareRobustCommand{\iconglobe}{{\color{blue!60!black}\faGlobe}}
\DeclareRobustCommand{\icondesktop}{{\color{gray!70!black}\faDesktop}}
\DeclareRobustCommand{\iconplay}{{\color{green!60!black}\scalebox{0.9}{\faIcon[regular]{play-circle}}}}
\DeclareRobustCommand{\iconstop}{{\color{red!60!black}\faStop}}
\DeclareRobustCommand{\iconnote}{{\color{blue!60!black}\faEdit}}
\DeclareRobustCommand{\iconlock}{{\color{gray!60!black}\faLock}}
\DeclareRobustCommand{\iconkey}{{\color{orange!70!black}\faKey}}
\DeclareRobustCommand{\iconfolder}{{\color{orange!60!black}\faFolderOpen}}
\DeclareRobustCommand{\iconrocket}{{\color{blue!60!black}\faRocket}}
\DeclareRobustCommand{\icongear}{{\color{gray!60!black}\faCog}}
\DeclareRobustCommand{\icontarget}{{\color{blue!60!black}\scalebox{0.9}{\faTasks}}}
\DeclareRobustCommand{\iconpc}{{\color{gray!70!black}\faLaptop}}
\DeclareRobustCommand{\iconcheck}{{\color{green!60!black}\faCheck}}
\DeclareRobustCommand{\iconpencil}{{\color{orange!70!black}\scalebox{0.9}{\faIcon{pencil-alt}}}}
\DeclareRobustCommand{\icontrophy}{{\color{orange!80!black}\scalebox{0.9}{\faTrophy}}}
\DeclareRobustCommand{\iconban}{{\color{red!70!black}\faBan}}
